\section{Una reconsideración de la clasificación L1-L5 de OpenAI}

\subsection{El significado de los cinco niveles}

OpenAI definió cinco niveles de capacidad de la IA: L1 ChatBot → L2 Reasoning → L3 Agent → L4 Innovation → L5 Organization.

\begin{importantbox}{Los niveles no guardan una relación lineal entre sí}
La intuición clave de Yang Zhilin (杨植麟) sobre L1-L5:
\begin{enumerate}
    \item Estos cinco niveles no son del todo secuenciales: algunas capacidades se desarrollan en paralelo
    \item El límite superior de Agent depende de la capacidad de Reasoning, pero Agent no tiene que esperar a que Reasoning madure por completo para empezar a desarrollarse
    \item El indicador de Innovation es que «el modelo participa en el desarrollo del modelo»; por ejemplo, K2 participa en el desarrollo de K3
    \item Ya ha aparecido un embrión de Organization: en los sistemas Multi-Agent, distintos Agent asumen distintos roles
\end{enumerate}
\end{importantbox}

\begin{knowledgebox}{Resolver los problemas de L3 con la tecnología de L4}
Hay una dependencia circular interesante: para resolver el problema de generalización de Agent (L3), puede hacer falta la tecnología de Innovation (L4), es decir, usar IA para entrenar y alinear IA. Esto muestra que entre estos niveles existe una relación de interacción compleja, y no una simple dependencia secuencial.
\end{knowledgebox}

\subsection{Resumen de la sección}

L1-L5 son varios hitos técnicos importantes, pero no guardan una relación estrictamente lineal. Reasoning y Agent son requisitos previos de Innovation y Organization, pero entre ellos existe una interdependencia compleja. La capacidad de cada nivel sigue mejorando de manera continua, sin un «tope» definido.
